\section{Обзор литературы}
\label{sec:related}

\subsection{Синтез медицинских изображений}

\paragraph{Эпоха классических методов.}
До широкого распространения глубоких порождающих моделей задача синтеза медицинских изображений решалась преимущественно средствами вычислительной геометрии, статистики формы и физики формирования изображения. Первая большая ветвь~--- цифровые анатомические фантомы. Обер-Броше и соавт.~\cite{aubert2006twenty} построили библиотеку из двадцати высокоразрешающих цифровых фантомов мозга на основе МРТ нормальных добровольцев из проекта BrainWeb. Каждый фантом представлял собой объёмную карту вероятностей принадлежности вокселя к одному из тканевых классов и использовался для генерации синтетических МРТ-данных с известным эталоном при разработке и валидации алгоритмов сегментации. Сегарс и соавт.~\cite{segars20104d} разработали четырёхмерный фантом XCAT (eXtended CArdiac-Torso), описывающий анатомию торса с помощью NURBS-поверхностей и моделирующий дыхательное и сердечное движения; XCAT поныне остаётся стандартом в задачах симуляции КТ- и ПЭТ-измерений и в задачах оптимизации протоколов лучевой дозы.

Вторая ветвь классики~--- статистические модели формы и текстуры. Активные модели внешнего вида (Active Appearance Models, AAM), предложенные Кутсом, Эдвардсом и Тэйлором~\cite{927467}, описывают объект совместным распределением координат опорных точек и интенсивностей в области, ограниченной этими точками, после устранения ковариации между формой и яркостью. AAM позволяли как параметрический синтез вариаций анатомии (путём сэмплирования из обученной модели в собственном базисе), так и анализ~--- регистрацию модели на новом изображении градиентным методом. Применения охватывали кардиальную МРТ, церебральную МРТ, дентальную рентгенографию и офтальмологию.

Все классические методы объединяет одно фундаментальное ограничение: распределение, из которого они порождают, лежит на низкоразмерном линейном многообразии в пространстве изображений (PCA-базис для AAM, заданное семейство фантомов для XCAT/BrainWeb). Это исключает воспроизведение тонкой патологической текстуры, межпациентной вариативности на уровне отдельных вокселей и нелинейных эффектов формирования изображения (артефакты, шум сканера, частичные объёмы). Кроме того, расширение модели на новую анатомическую область требует ручной разметки сотен исследований.

\paragraph{Глубокие порождающие модели.}
Качественный сдвиг произошёл с распространением глубоких порождающих моделей. Первое поколение опиралось на генеративно-состязательные сети~(GAN)~\cite{goodfellow2014generative}: два состязающихся параметризованных распределения~--- порождающее и распознающее~--- сходятся к равновесию, при котором сгенерированные образцы статистически неотличимы от реальных. В задаче трансляции МРТ в КТ Ху и соавт.~\cite{hu2021bidirectional} предложили двунаправленную GAN-архитектуру, обеспечивающую взаимно согласованные преобразования. Для синтеза ПЭТ из МРТ разработаны специализированные фреймворки на основе условных GAN~\cite{WEI2019101546, ZHANG2022106676, WANG2024102983}, нередко с трёхмерными сверточными кодировщиками~\cite{WANG2024102983} и механизмами self-attention~\cite{https://doi.org/10.1002/mrm.28819}. Для задачи суперрезолюции предложены многомасштабные сети~\cite{pham2019multiscale} и GAN-CIRCLE~\cite{you2019ct}. Отдельного упоминания заслуживают работы по синтезу данных для реконструкции~\cite{peng2022towards, xie2022measurement}, где диффузионные модели позволяют решать проблему неполных измерений. Несмотря на впечатляющие результаты, GAN-подход имеет хорошо известные недостатки: коллапс мод, нестабильность обучения и отсутствие согласованной меры правдоподобия, что затрудняет их использование в задачах с требованием точной воспроизводимости числовых значений~--- таких как КТ в HU.

Диффузионные вероятностные модели~\cite{ho2020denoising, sohl2015deep} принесли очередной качественный скачок. Их теоретическое преимущество над GAN~--- стабильное обучение с явным правдоподобием, а недостаток~--- медленный многошаговый инференс. Ромбах и соавт.~\cite{rombach2022high} показали, что латентные диффузионные модели способны синтезировать изображения высокого разрешения при умеренных вычислительных затратах. В медицинской визуализации они нашли применение для гармонизации МРТ~\cite{durrer2023diffusion}, перевода МРТ в КТ~\cite{graf2023denoising}, трёхмерного синтеза~\cite{guo2025maisi, zhu2023make} и цикловых MRI-КТ задач~\cite{pan2025cycle}. Тем не менее ключевым ограничением диффузионных методов остаётся их медлительность: получение одного образца требует нескольких сотен шагов денойзинга, что неприемлемо для клинических приложений реального времени и для обучения downstream-моделей на синтетических данных.

\subsection{Трансляция КТ-контраста}

Задача синтеза виртуальных нативных (или контрастных) КТ-изображений рассматривалась ещё в контексте двухэнергетической КТ~(ДЭКТ). Вирар и соавт.~\cite{VIRARKAR2022293} дают обширный клинический обзор приложений ДЭКТ для виртуального нативного изображения в брюшной полости и малом тазу. Кин и соавт.~\cite{article_dual} предложили генерацию глубокого виртуального нативного КТ с применением двухслойной ДЭКТ для радиотерапевтического планирования. Пан и соавт.~\cite{pan2024synthetic} и Каплантар и соавт.~\cite{article_nature} использовали CycleGAN для синтеза нативной серии из контрастных МРТ и КТ соответственно, обходясь без парных обучающих данных. Ким и соавт.~\cite{kim2024adaptive} разработали адаптивную латентную диффузионную модель для трёхмерного межмодального МРТ-синтеза.

Наиболее близкой к предлагаемому методу является работа Лю и соавт.~\cite{liu2025see}~--- метод SMILE, применяющий анатомически-управляемый диффузионный подход к задаче контрастного усиления. Авторы вводят анатомическое предобусловливание и специальные потери, однако базовая архитектура DDPM сохраняет многошаговый инференс. На наших данных SMILE показал значительно более высокую ошибку MAE (30.22\,HU против 4.25\,HU у MFM), что, по всей видимости, связано со сдвигом распределений между данными, использованными для создания предобученной модели, и нашим датасетом.

\subsection{Математические предпосылки: оптимальный транспорт и непрерывные потоки}

Флоу матчинг опирается на классическую теорию переноса меры. Монж в 1781~году поставил задачу о минимизации суммарной стоимости перевозки массы; Канторович~\cite{kantorovich1942translocation} обобщил её до релаксации, допускающей расщепление массы. Бренье~\cite{brenier1991polar} доказал существование и единственность оптимального транспортного отображения при квадратичной стоимости и абсолютно непрерывной начальной мере; в дальнейшем теория получила систематическое изложение в работах Виллани~\cite{villani2008optimal} и Сантамброджо. Геометрическая динамическая переформулировка~--- задача Бенаму--Бренье~\cite{benamou2000computational}~--- сводит расстояние Вассерштейна $W_2^2(\pi_0,\pi_1)$ к минимизации кинетической энергии непрерывного потока, удовлетворяющего уравнению непрерывности. Именно эта переформулировка лежит в основе современного флоу матчинга.

Параллельная линия исследований развивает порождающие модели на основе непрерывных нормализующих потоков~--- Neural ODE~\cite{chen2018neural} и FFJORD~\cite{grathwohl2018ffjord}: распределение задаётся как образ простой меры через решение ОДУ с обучаемой правой частью. Главное препятствие исторически~--- необходимость симулировать траекторию и считать определитель якобиана при обучении максимума правдоподобия.

\subsection{Флоу матчинг}

Симуляционно-свободная парадигма обучения непрерывных потоков предложена практически одновременно Липманом и соавт.~\cite{lipman2022flow}, Алберго и Ванден-Ейнденом~\cite{albergo2022building} (формализм стохастических интерполянтов) и Лю и соавт.~\cite{liu2022flow} (Rectified Flow). Идея состоит в том, чтобы соединить шумовое распределение с распределением данных посредством явного интерполянта $\mathbf{x}_t=\alpha_t\mathbf{x}_0+\sigma_t\mathbf{x}_1$ и обучить нейронную сеть предсказывать условную производную $\partial_t\mathbf{x}_t$; теорема о градиентной эквивалентности гарантирует, что такая регрессия минимизирует расхождение между сетью и истинным маргинальным полем скоростей. Тонг и соавт.~\cite{tong2023improving} обобщили подход с использованием минибатчного оптимального транспорта, получив более прямые пути и тем самым улучшив качество одношагового инференса. Сонг и соавт.~\cite{song2020score} ранее показали, что само диффузионное обучение можно интерпретировать как обучение поля скоростей вероятностного потока ОДУ.

В медицинской визуализации флоу матчинг только начинает завоёвывать позиции. Яздани и соавт.~\cite{yazdani2025flow} рассмотрели его для синтеза медицинских изображений в целом, продемонстрировав баланс между скоростью и качеством. Рейно и соавт.~\cite{reynaud2025echoflow} разработали EchoFlow~--- фундаментальную модель для генерации ультразвуковых изображений сердца. При этом ни одна из известных работ не применяла флоу матчинг к задаче парной трансляции КТ-контраста и не предлагала специализированную архитектуру для этого класса задач.

\subsection{Базовые архитектуры}

Значительная часть предлагаемых экспериментальных сравнений опирается на архитектуры, хорошо зарекомендовавшие себя в медицинской сегментации. U-Net~\cite{ronneberger2015u} --- классическая архитектура с энкодером-декодером и пропускными связями --- остаётся стандартом де-факто для задач с попиксельным предсказанием. SegResNet~\cite{myronenko20183d} расширяет эту схему с помощью остаточных блоков и вариационного автокодировщика для регуляризации. SwinUNETR~\cite{hatamizadeh2021swin} заменяет сверточный энкодер иерархическим трансформером на основе sliding window Swin, что позволяет захватывать долгодистанционные зависимости при умеренной вычислительной стоимости. UMambaBot~\cite{ma2024u} включает модели пространства состояний~(Mamba) в U-Net фреймворк для эффективного моделирования длинных зависимостей в биомедицинских изображениях.

Для сравнения с диффузионными моделями используются DDPM~\cite{ho2020denoising} с условной подачей исходного изображения и ускоренным семплированием DDIM~\cite{song2020denoising}. Downstream-сегментация оценивается nnUNet\,v2~\cite{isensee2024nnu}~--- самоконфигурирующимся фреймворком, задающим стандарт в данной области.
